\chapter{Verbs}
\label{chap:verbs}

\section{Verb forms}

\subsection{Conjugation}

Each verb has both finite and non-finite forms. Finite forms are used to express tense, aspect, mood, and person, while non-finite forms are used mainly as nominal or adjectival elements in a sentence.

Conjugation in Niklorsk is minimal with verbs only inflecting for tense, aspect and mood. There are two main conjugation classes, depending on whether the verb ends in a consonant or a vowel. The verbs \nikp{hjalt}{to write} and \nikp{bra}{to give} are used as examples below to illustrate the conjugation of verbs in Niklorsk.

Inflectional endings are stress-neutral and preserve the lexical stress of the verb (see Chapter \ref{chap:phonology}, \S\ref{sec:stress-prosody}). Thus \nk{ˈhjalt} retains initial stress in \nk{ˈhjaltar} and \nk{ˈhjaltatt}. These spellings show stress only; the marks are not part of ordinary orthography.

\begin{table}[!ht]
    \small
    \begin{tblr}{
        colspec={XX[1.75]X[1.75]}
    }
        {\scshape finite forms} &\\
        & {\scshape class i} & {\scshape class ii}\\
        Present & \nk{hjalt} & \nk{bra}\\
        Progressive & \nk{hjaltatt} & \nk{bratt}\\
        Past & \nk{hjaltar} & \nk{brar}\\
        Perfect & \nk{hjaltet} & \nk{bratt}\\
        Future & \nk{hjaltäs} & \nk{brask}\\
        Conditional & \nk{hjaltårs} & \nk{brars}\\
        Neg. Conditional & \nk{hjalter} & \nk{brager}\\
        & \\
        {\scshape non-finite forms} &\\
        Infinitive & \nk{hjalts} & \nk{brags}\\
        Perf. Infinitive & \nk{hjaltsk} & \nk{bragsk}\\
        Present participle & \nk{hjaltade} & \nk{brade}\\
        Past participle & \nk{hjaltagne} & \nk{bragne}\\
        Imperative & \nk{hjalte} & \nk{brage}\\
    \end{tblr}
    \caption{Conjugation template for Niklorsk verbs}
    \label{tab:verb-conjugation}
\end{table}

After Class I vowels ending in \nk{-s}, an epenthetic \nk{-e-} is inserted before the infinitive suffix \nk{-s}, e. g., \nikp{bauses}{to kill} from \nikp{baus}{to kill}.

The participial endings \nk{-ade} and \nk{-agne} trigger a change in the vowel of native Class I monosyllabic verbs. In particular, the following changes occur: (1) \nk{æ} $\rightarrow$ \nk{ö}, (2) \nk{au} $\rightarrow$ \nk{o}, (3) \nk{o} $\rightarrow$ \nk{u}, (4) \nk{å} $\rightarrow$ \nk{a} and (5) \nk{ee} $\rightarrow$ \nk{o}. Examples are given in Table \ref{tab:vowel-changes-participles} below.

\begin{table}[!ht]
    \small
    \begin{tblr}{
        colspec = {X[1.5]XXXXX}
    }
    & \nk{væd} & \nk{naud} & \nk{sond} & \nk{mårs} & \nk{leet}\\
    & {\footnotesize `to see'} & {\footnotesize `to enter'} & {\footnotesize `to speak'} & {\footnotesize `to dwell'} & {\footnotesize `to walk'}\\
    Pres. Part. & \nk{vödade} & \nk{nodade} & \nk{sundade} & \nk{märsade} & \nk{lotade}\\
    Past Part. & \nk{vödagne} & \nk{nodagne} & \nk{sundagne} & \nk{märsagne} & \nk{lotagne}\\
    \end{tblr}
    \caption{Vowel changes in the present and past participles of native Class I monosyllabic verbs}
    \label{tab:vowel-changes-participles}
\end{table}

The present tense form is the base form of the verb, and is used as the verb's lemma or dictionary form. Most verbs follow the conjugation pattern of \nikp{hjalt}{to write} given above. However, there are a few irregular verbs, most notably the verbs \nikp{ö}{to be} (whose conjugation is given in Table \ref{tab:verb-to-be}), \nikp{vor}{there is} (v. Table \ref{tab:verb-there-is}) and \nikp{lod}{to have}.

Latinate loanwords often take the ending \nk{-er} and are conjugated as follows, with the verb \nikp{kanseler}{to cancel} as an example.

\begin{table}[!ht]
    \small
    \begin{tblr}{
        colspec = {XX}
    }
        {\scshape finite forms} &\\
        Present & \nikp{kanseler}{(I) cancel}\\
        Progressive & \nikp{kanselatt}{(I) am cancelling}\\
        Past & \nikp{kanselera}{(I) cancelled}\\
        Perfect & \nikp{kanselert}{(I) have cancelled}\\
        Future & \nikp{kanselisk}{(I) will cancel}\\
        Conditional & \nikp{kanselerårs}{(I) would cancel}\\
        Neg. Conditional & \nikp{kanselerer}{(I) would not cancel}\\
        & \\
        {\scshape non-finite forms} &\\
        Infinitive & \nikp{kanselers}{to cancel}\\
        Present participle & \nikp{kanselade}{cancelling}\\
        Past participle & \nikp{kanselagne}{cancelled}\\
        Imperative & \nikp{kansel}{cancel!}\\
    \end{tblr}
    \caption{Conjugation of the verb \nikp{kanseler}{to cancel}}
    \label{tab:verb-conjugation-latinate}
\end{table}

Verbs formed with \nk{-iser/-tiser} follow this pattern, with primary stress on the syllable containing the suffix's first vowel (see Chapter \ref{chap:derivations}, \S\ref{sec:suffix-iser}). Replacing or deleting final \nk{-er} leaves that stressed vowel in place, as in \nk{normaˈliser} $\rightarrow$ \nk{normaˈlisatt}. The action noun \nk{normaˈliserats} retains the same stress because derivational \nk{-ats} is stress-neutral.

\section{Finite forms}

\subsection{Present and progressive tenses}

The present tense is used to describe actions that are currently happening or that happen regularly. The progressive tense\footnote{
    The progressive `tense' is more properly analyzed as an aspect, as it does not indicate time but rather the ongoing nature of an action. However, for simplicity, and following the traditional terminology, we will refer to it as a tense.
} is used to describe actions that are currently ongoing (i.e., at the time of speaking) or that was/will be ongoing at a specific time in the past or future. The present tense uses the base form of the verb, while the progressive tense uses the base form of the verb followed by the suffix \nk{-att} or \nk{-tt} for Class I and Class II verbs, respectively. The present tense may be used to describe an ongoing event in the present, but if the emphasis is on the ongoing nature of the event, the progressive tense is preferred.

\pex
\begingl
\gla Ngöld börd bot. En ngöld bördatt bot.//
\glb cow eat.\Pres{} grass the.\M{} cow eat.\Prog{} grass//
\glft `Cows eat grass (in general). The cow is eating grass (right now).'//
\endgl
\xe

\pex
\begingl
\gla I kvett slämatt nar de hus.//
\glb in previous:day paint.\Prog{} \Third{}\Sg{} the.\F{} house//
\glft `Yesterday he was painting the house.'//
\endgl
\xe

The present tense is used for general truths or habitual actions. It is thus used in proverbs, instructions or recipes, in scientific or mathematical statements, and in announcements (often in the passive voice). Where the progressive tense is used, it refers to a concrete event or action.

\pex
\begingl
\gla Eg magasin spören i måment.//
\glb the.\N{} store close.\Pres{}.\Pass{} in moment//
\glft `The store is closing in a moment.'//
\endgl
\xe

\pex
\begingl
\gla Cykliske hrådd ö hrådd, kan generareren på jörv hamde element.//
\glb cyclic.\M{} group be.\Pres{} group, what generate.\Pres{}.\Pass{} by single.\M{} \Refl{} element//
\glft `A cyclic group is a group that is generated by a single element.'//
\endgl
\xe

The present tense may also be used in journalistic reporting (particularly in headlines) or in literary narratives for stylistic reasons, even when the even has already occurred and is yet to occur.

\pex
\begingl
\gla Går kom dänn krust i kvett.//
\glb ten die.\Pres{} in earthquake in previous:day//
\glft `Ten die in yesterday's earthquake.'//
\endgl
\xe

In spoken Niklorsk, preference is usually given to the present tense over the future tense when referring to events in the near future.

\subsection{Past and perfect tenses}

The past tense is used to describe actions that have occured in the past. The perfect tense is used to describe actions that have been completed at the time of speaking or at a specific time in the past or future.


\subsection{Negation of verbs}

Verbs can be negated by the particle \nk{såm}, placed directly after the main verb in the sentence (i.e., in the case of modal verbs, after the modal verb). 

\pex
\begingl
\gla Nar börd såm skäll.//
\glb \Third{}\Sg{} eat not fish.//
\glft `He does not eat fish.'//
\endgl
\xe

\pex
\begingl
\gla Lessi vell såm dansers.//
\glb Lessi can not dance.\Inf{}//
\glft `Lessi cannot dance.'//
\endgl
\xe


\subsection{Conditional form}

The conditional form (sometimes referred to as the subjunctive mood in some texts) is used mainly for hypothetical, unreal, or speculative statements. 

In conditional statements, the conjunctions \nikp{efte}{if} and \nikp{efte bär}{if ever} triggers the use of the conditional form in the conditional clause. The conjunction \nikp{efte ngå}{if only} triggers the use of the conditional form in both the conditional clause and the result clause.

\pex
\begingl
\gla Sortäs ag hamdär nyst, efte sortårs Matti eg hamde.//
\glb complete.\Fut{} \First{}\Sg{} \Refl{}.\F{} part, if complete.\Cond{} Matti the.\Nt{} \Refl{}.\N{}//
\glft `I will complete my part if Matti completes his.'//
\endgl
\xe

Some specific constructions require the conditional form, including, among others, the following:

\begin{enumerate}[wide, nosep, label=(\alph*)]
    \item With the adverb \nikp{nå}{such, so} and an adjective or adverb, to express a hypothetical or unreal situation, or the result of the action in the main clause.
    \pex
    \begingl
    \gla Ag ö nå maus, megd togårs ag stetts.//
    \glb \First{}\Sg{} be.\Pres{} so happy, that want.\Cond{} \First{}\Sg{} cry.\Inf{}//
    \glft `I am so happy I want to cry.'//
    \endgl
    \xe
    \pex
    \begingl
    \gla Nar sondar nå rontæg, megd fallårs ag när rings.//
    \glb \Third{}\Sg{} speak.\Pst{} so fast, that cannot.\Cond{} \First{}\Sg{} him understand.\Inf{}//
    \glft He spoke so fast I wasn't able to understand him.//
    \endgl
    \xe
    \item With superlative conjunctions, equivalent to English `the \textit{most} (adjective) \textit{that} (verb)'.
    \pex
    \begingl
    \gla Hon ö de strålinner, käll vædårs ag nörs.//
    \glb she be.\Pres{} the.\F{} pretty.\Comp{}.\F{} who see.\Cond{} \First{}\Sg{} never//
    \glft `She is the prettiest (girl) I have ever seen.'//
    \endgl
    \xe
\end{enumerate}

\subsection{Modal verbs}
\label{sec:modal-verbs}

Modal verbs appear before the main verb in a sentence and are used to express necessity, possibility, ability, or permission. They are followed by the infinitive form of the main verb. Then infinitive appears directly after the modal verb, except when the main verbs takes a pronoun or demonstrative as its object, in which case the infinitive appears after the object.

\pex
\begingl
\gla Andärs fall væds Matti i krugnat. Andärs fall när væds.//
\glb Andärs cannot see.\Inf{} Matti in crowd Andärs cannot \Third{}\Sg{}.\Acc{} see.\Inf{}//
\glft `Andärs cannot see Matti in the crowd. Andärs cannot see him.'//
\endgl
\xe

\pex
\begingl
\gla Ag tog såm med mykks.//
\glb \First{}\Sg{} want.\Pres{} not that do.\Inf{}//
\glft `I don't want to do that.'//
\endgl
\xe

Adverbs (except those of time and place formed using a preposition) may optionally appear between the modal verb and the main verb. Where the main verb takes a pronouns or demonstrative as its object, the adverb must always appear after the object.

\pex
\begingl
\gla Meg dägn dag begns en afte. Meg dägn när dag begns.//
\glb \Second{}\Sg{} must again read.\Inf{} the.\Nt{} book \Second{}\Sg{} must \Third{}\Sg{}.\Acc{} again read.\Inf{}//
\glft `You must read the book again. You must read it again.'//
\endgl
\xe


Ability is expressed by the modal verbs \nk{vell} and \nk{igom}. \nk{Vell} is used to express general ability, particularly those that are learned or acquired, while \nk{igom} is used to express physical ability, or the ability to do perform an action without obstruction. Both \nk{vell} and \nk{igom} have separate negative forms, \nk{kun} and \nk{fall}, respectively. The aforementioned modal verbs conjugate regularly except for \nk{igom} which has the irregular past form \nk{ignar}.

\pex
\begingl
\gla Ag vell sonds Engelsk.//
\glb \First{}\Sg{} can speak.\Inf{} English.//
\glft `I can speak English.'//
\endgl
\xe
\pex
\begingl
\gla Nar igom malms atta magasin, efte leker Martin saugn.//
\glb \Third{}\Sg{} can go.\Inf{} to:the store, if respond.\N{}\Cond{} Martin us//
\glft `He can go to the store if Martin doesn't respond to us.'//
\endgl
\xe

The construction \nikp{Med igom, meg\dots}{It can be that\dots} and its negative counterpart \nikp{Med fall, meg\dots}{It cannot be that\dots} can be used to express doubt or uncertainty in addition to possibility, similar to the English construction `It may be that\dots' or `It is possible that\dots'. The verb in the subordinate clause takes the conditional form.

\pex
\begingl
\gla Med igom, meg igomer ag mosts atta fjeltt.//
\glb that can.\Pres{} \Rel{} can.\N{}\Cond{} \First{}\Sg{} come.\Inf{} to:the party//
\glft `It's possible that I wouldn't be able to come to the party.'//
\endgl
\xe

Necessity is expressed by \nikp{dägn}{must} or \nikp{mått}{should}. \nk{Dägn} is used to express a strong necessity or obligation, while \nk{mått} often takes a more suggestive or advisory tone.

\pex
\begingl
\gla Ag malm i spa; meg mått aka provers.//
\glb \First{}\Sg{} go.\Pres{} to spa \Second{}\Sg{} should also try.\Inf{}//
\glft `I'm going to the spa; you should also try it.'//
\endgl
\xe

\pex
\begingl
\gla Meg dägn såm hents at mama saugna börn.//
\glb \Second{}\Sg{} must not tell.\Inf{} to mother our secret//
\glft `You must not tell mother our secret.'//
\endgl
\xe

\section{Non-finite forms}

\subsection{The infinitive}


The infinitive is used in certain compound verb constructions as the object of the main verb. It is used with the modal auxiliaries \nikp{vell}{can}, \nikp{igom}{can}, \nikp{kun}{cannot}, \nikp{fall}{cannot}, \nikp{tog}{want}, \nikp{dägn}{must}, \nikp{mått}{should}, etc. and verbs with quasi-modal interpretation like \nikp{geed}{hope}, \nikp{kvill}{begin}, \nikp{kunda}{stop} etc. Similarly, the infinitive is used as the object of sensory verbs like \nikp{væd}{see,} \nikp{begn}{to hear}, etc.

\pex
\begingl
\gla Nar tog malms at Norge.//
\glb \Third{}\Sg{} want.\Pres{} go.\Inf{} to Norway.//
\glft `He wants to go to Norway.'//
\endgl
\xe

\pex
\begingl
\gla Matti dägn myks med smed leeb.//
\glb Matti must do.\Inf{} that before morrow.//
\glft `Matti must do that before tomorrow.'//
\endgl
\xe

\pex
\begingl
\gla Ag kvillar sonds.//
\glb \First{}\Sg{} begin.\Pres{} speak.\Inf{}//
\glft `I began to speak.'//
\endgl
\xe

In secondary clauses where the subject of the infinitive is the object of the main verb, the infinitive is used as the preferred alternative to a complement clause.

\pex
\begingl
\gla Nar vædar Matti kupps en afte. Nar vaedar, meg kuppar Matti en afte.//
\glb \Third{}\Sg{} see.\Pst{} Matti take.\Inf{} the.\Nt{} book. \Third{}\Sg{} see.\Pst{}, that take.\Pst{} Matti the.\Nt{} book.//
\glft `He saw Matti take the book. He saw that Matti took the book.'//
\endgl
\xe

The perfect infinitive is used in a similar manner to the infinitive, but with the action denoted by the verb taking a perfective interpretation; i.e., the action is completed or has been completed at the time referred to by the main verb. The perfect infinitive is formed by adding the suffix \nk{-(i)sk} to the infinitive form of class I verbs and \nk{-gsk} to the infinitive form of class II verbs.

\pex
\begingl
\gla Rett är rammsk av mæg i man projekt.//
\glb pleasure be.\Pst{} work.\Perf.\Inf{} with you.\Acc{} in this project//
\glft `It was a pleasure to have worked with you on this project.'//
\endgl
\xe

\pex
\begingl
\gla Matti tind fjerdsk hamde mobil.//
\glb Matti seems.\Pres{} forget.\Perf.\Inf{} \Refl{}.\M{} phone//
\glft `Matti seems to have forgotten his phone.'//
\endgl
\xe


The infinitive is used with the preposition \nikp{i}{in, on, at} to indicate an action that is about to take place or is in the process of taking place at the same time as the action of the main verb. This construction correponds to English `on + gerund,' `upon + gerund,' or similar gerundial phrases. The \nk{i} + infinitive clause can also take a causative reading, corresponding to English `because of doing something,' etc. The \nk{i} + infinitive clause should have the same subject as the main clause, and is usually fronted in the sentence.

\pex
\begingl
\gla I föds en kaffe klörar ag hæm.//
\glb in heat:up.\Inf{} the coffee burn.\Pst{} \First{}\Sg{} \Refl{}//
\glft `While heating up the coffee, I burned myself.'//
\endgl
\xe

\pex
\begingl
\gla I begns de dokument kostar nar en dreg.//
\glb in read.\Inf{} the.\F{} document, change.\Pst{} \Third{}\Sg{} the.\M{} mind//
\glft `Upon reading the document, he changed his mind.'//
\endgl
\xe

\pex
\begingl
\gla Martin fall megers i lods som 18 spisar.//
\glb Martin cannot vote.\Inf{} in have.\Inf{} \Neg{} 18 year.\Pl{}//
\glft `Martin cannot.\Pres{} vote, being under 18 years old.'//
\endgl
\xe

\pex
\begingl
\gla Andärs striggar hæm i sæds talt.//
\glb Andärs faint.\Pst{} \Refl{}.\Acc{} in hear.\Inf{} news//
\glft `Anders fainted upon hearing/because of the news.'//
\endgl
\xe

To emphasize the causative meaning, the preposition \nikp{vin}{for, because of} is used instead of \nk{i}. This is especially common in phrases like \nikp{mæga vin...}{thanks for...} or \nikp{hjörtat vin...}{sorry for...}.

\pex
\begingl
\gla Mæga vin hents ægn det med.//
\glb thanks for tell.\Inf{} \First{}\Sg{}.\Acc{} about that//
\glft `Thanks for telling me about that.'//
\endgl
\xe


\pex
\begingl
\gla Hjörtat vin opslagsk såm ja billä.//
\glb sorry for send.\Perf.\Inf{} not it.\Acc{} right:away//
\glft `Sorry for not having sent it right away.'//
\endgl
\xe

The infinitive is used after sequential conjunctions such as \nikp{smed}{before} and \nikp{åd}{after}, where English would use a gerund.

\pex
\begingl
\gla Ag svorar åd begns en afte.//
\glb \First{}\Sg{} sleep.\Pst{} after read.\Inf{} the.\Nt{} book//
\glft `I slept after reading the book.'//
\endgl
\xe

\pex
\begingl
\gla Greede som en mobil æder vands.//
\glb use.\Imp{} not the.\M{} phone while drive.\Inf{}//
\glft `Don't use the phone while driving.'//
\endgl
\xe

In colloquial speech, as well as in written signages, the infinitive is used for prohibitions. For examples, \nikp{Rauks som.}{No smoking.} or \nikp{Pjæts som.}{No loitering.}

\section{Separable verbs}

Separable verbs or coverbs (Niklorsk \nikp{koverb}{coverb}) are verbs that consist of a verb and a preposition or adverb that modifies the meaning of a verb. By default, the preposition or adverb appears before the verb, with the verb spelled as a single word.

\pex
\begingl
\gla Ag anslötar sla hamde brord vin hjekks av Matti.//
\glb \First{}\Sg{} come:out.\Pst{} from \Refl{}.\M{} room in:order:to talk.\Inf{} with Matti//
\glft `I came out of my room to talk to Matti.'//
\endgl
\xe

Common prepositions used in separable verbs include \nikp{an}{out}, \nikp{av}{with}, \nikp{i}{in}, \nikp{at}{into}, \nikp{hel}{up}, \nikp{hæt}{down}, \nikp{een}{off, away}, \nikp{att}{together}, \nikp{op}{apart}, \nikp{be}{perfective}.

The addition of a prefix to the verb changes the meaning of the verb. With most verbs, especially verbs of motion, the prefix often takes a literal meaning. For example, from \nikp{slöt}{come} we get \nikp{anslöt}{come out}, \nikp{eenslöt}{come off}, \nikp{atslöt}{come in}. Some separable verbs, however, take a figurative/idiomatic meaning, e.g., \nikp{hjalt}{to write} but \nikp{eenhjalt}{to write off (as a loss)}, and \nikp{betug}{to become fat} but \nikp{eentug}{to fatten away, to become obese}. Some separable verbs may take both literal and idiomatic meanings, as in the above example with the verb \nikp{hætslöt}{to come down}:

\pex
\begingl
\gla En ängel hætslötar sla bærg.//
\glb the.\M{} angel come:down.\Pst{} from heaven//
\glft `The angel came down from heaven.'//
\endgl
\xe

\pex
\begingl
\gla Sä pris sla veetöt slötar såm hæt.//
\glb the.\Pl{} price of wheat.\Gen{} come.\Pst{} not down//
\glft `The price of wheat did not come down.'//
\endgl
\xe


The preposition or adverb is split from the verb in the following cases:
\begin{enumerate}[wide, nosep, label=(\alph*)]
    \item When the verb is negated, e.g., by the particle \nikp{såm}{not}, in which case the preposition or adverb is placed after the negation particle.
    
    \pex
    \begingl
    \gla Ag slötar såm an sla hamde brord.//
    \glb \First{}\Sg{} come.\Pst{} not out from \Refl{}.\M{} room//
    \glft `I did not come out of my room.'//
    \endgl
    \xe

    \item When the verb is used in the passive voice.
    
    \pex
    \begingl
    \gla En samm braren hel par brevet.//
    \glb the.\M{} land give.\Pst{} up by:the owner//
    \glft `The land was given up by the owner.'//
    \endgl
    \xe

    \item When the verb takes a pronominal or demonstrative object, in which case the preposition or adverb is placed after the object.
    
    \pex
    \begingl
    \gla Ag hjaltar med een för operationskost.//
    \glb \First{}\Sg{} write.\Pst{} that off as operating:cost//
    \glft `I wrote that off as an operating cost.'//
    \endgl
    \xe

    \item In questions, where the preposition or adverb is placed after the verb, subject to other rules of question formation.
    
    \item With modal and quasi-modal verbs, where the prefix appears before the modal or quasi-modal verb.
    
    \item In imperatives.
\end{enumerate}

\section{Copular and existential constructions}

\subsection{The verbs \nk{ö} and \nk{går}}

The verb \nikp{ö}{to be} and \nikp{går}{to become} serve as the copulae in Niklorsk. It is used to link the subject of a sentence to a subject complement (discussed in detail below). The verb \nikp{ö}{to be} is slightly irregular and has the following conjugation:

\begin{table}[!ht]
    \small
    \begin{tblr}
        {
            colspec={X[1.25,l]X[1.5,l]X[1.75,l]}
        }
        Finite forms & & \\
        Present & \nikp{ö}{I am} & \nikp{går}{I become}\\
        Progressive & \nikp{ett}{I am being} & \nikp{gät}{I am becoming}\\
        Past & \nikp{är}{I was} & \nikp{är}{I became}\\
        Perfect & \nikp{ätt}{I have been} & \nikp{gåt}{I have become}\\
        Future & \nikp{isk}{I will be} & \nikp{gikt}{I will become}\\
        Conditional & \nikp{års}{I would be} & \nikp{gins}{I would become}\\
        Negative Conditional & \nikp{av}{I would not be} & \nikp{gav}{I would not become}\\
        & & \\
        Non-finite forms & & \\
        Infinitive & \nikp{ös}{to be} & \nikp{gårs}{to become}\\
        Present participle & \nikp{edde}{being} & \nikp{gäde}{becoming}\\
        Past participle & \nikp{ägne}{been} & \nikp{gånge}{become}
    \end{tblr}
    \caption{Conjugation of the verb \nikp{ö}{to be}}
    \label{tab:verb-to-be}
\end{table}

\nk{Går} is used to indicate a change in state or condition, while \nk{ö} is used to indicate a permanent or inherent state or condition. The future tense form \nikp{isk}{I will be} is almost never used in spoken Niklorsk, with the future tense form of \nk{går} replacing it in most contexts.

\pex
\begingl
\gla Ag gikt unga, efte sittårs meg lære.//
\glb \First{}\Sg{} become.\Fut{} sad if leave.\Cond{} \Second{}\Sg{} now//
\glft `I will be sad if you leave now.'//
\endgl
\xe


\subsection{The existential verb \nk{vor}}

The verb \nikp{vor}{there is} is used to indicate the existence of something. It is also used as a copula, in contrast with \nikp{ö}{to be}, to describe a physical or temporary state or condition. The verb \nikp{vor}{there is} is conjugated as follows:

\begin{table}
    \small
    \begin{tblr}{
        colspec={XXX}
    }
        Finite forms & Present & \nikp{vor}{there is}\\
        & Progressive & \nikp{voratt}{there is being}\\
        & Past & \nikp{är}{there was}\\
        & Perfect & \nikp{vorst}{there has been}\\
        & Future & \nikp{voräs}{there will be}\\
        & Conditional & \nikp{vorems}{there would be}\\
        & Negative Conditional & \nikp{voren}{there would not be}\\
        & & \\
        Non-finite forms & Infinitive & \nikp{vors}{to be}\\
        & Present participle & \nikp{vorade}{being}\\
        & Past participle & \nikp{voragne}{been}
    \end{tblr}
    \caption{Conjugation of the existential verb \nikp{vor}{there is}}
    \label{tab:verb-there-is}
\end{table}

\nk{Vor} not \nk{ö} is used when the predicate refers to the location of the subject.

\pex
\begingl
\gla Ga vor dänn hus. Ag vor aka mer.//
\glb \First{}\Pl{} be in house. \First{}\Sg{} be also here.//
\glft `We are in the house. I am here.'//
\endgl
\xe

Thus, in the above example, both \nk{Ga ö dänn hus} and \nk{Ag ö aka mer} would be ungrammatical.

\section{The passive voice}

The passive voice in Niklorsk may be formed in one of two ways. The simple passive is formed by adding the suffix \nk{-(e)n} at the end of the conjugated verb form (in this case, the final /n/ of the suffix is dropped in the pronunciation). The subject of the active voice sentence becomes the object of the passive voice sentence, and the agent (if mentioned) is introduced by the preposition \nikp{på}{by} or \nikp{par}{by the}. Alternatively, the passive may also be formed by using the appropriate conjugation of the verb \nikp{vor}{there is} followed by the past participle of the main verb. The participle must then agree in gender and number with the subject of the sentence.

\pex
\begingl
\gla Matti kvintar eg hjålt.//
\glb Matti bake.\Pst{} the.\Nt{} cake//
\glft `Matti baked the cake.'//
\endgl
\xe

\pex
\begingl
\gla Eg hjålt dägn kvintsen i sægge.//
\glb the.\Nt{} cake must bake.\Inf{}.\Pass{} in morning//
\glft `The cake must be baked today.'//
\endgl
\xe

\pex
\begingl
\gla Eg hjålt är kvintagne på Matti.//
\glb the.\Nt{} cake be.\Pst{} bake.\PstP{}.\Nt{} by Matti//
\glft `The cake was baked by Matti.'//
\endgl
\xe

The \nk{vor} passive is often used in reporting an event that is isolated or unexpected, often with an adverbial phrase indicating the time or place of the event. It is also preferred where the speaker wants to emphasize the resulting state of the subject.

\pex
\begingl
\gla De jårt är strikkagnär par velt.//
\glb the.\F{} town be.\Pst{}. destroy.\PstP{}.\F{} by:the storm//
\glft `The town was destroyed by the storm.'//
\endgl
\xe

\pex
\begingl
\gla En bliv är daussagne på auto i kvett.//
\glb the.\M{} man be.\Pst{} hit.\PstP{}.\M{} by car in previous:day//
\glft `The man was hit by a car yesterday.'//
\endgl
\xe

In contrast, the simple passive is preferred when the action is not restricted to a particular time, as when speaking about actions that are habitual or repeated or which express general norms or truths. Instructions, announcements, and other impersonal statements are also often expressed in the simple passive and never in the \nk{vor} passive. This also contrasts with English, where the imperative is often used in the above contexts.


\pex
\begingl
\gla Sä patator mursten i kelv.//
\glb the.\Pl{} potato.\Pl{} cut.\Pres{}.\Pass{} in half//
\glft `The potatoes are cut in half.' or 'Cut the potatoes in half.'//
\endgl
\xe

\pex
\begingl
\gla Hamder hål hjalten mer.//
\glb \Refl.\F{} name write.\Pres{}.\Pass{} here//
\glft `One's name is written here.'//
\endgl
\xe

With modal verbs, the simple passive is used.

% personal vs impersonal
